% GENERATED FILE. Edit canonical lessons, then run npm run generate:books.
% canonical-source-hash: fnv1a64:52377bc17803fa8a
% canonical-lessons: KA-C217-agga, KA-C217-khali, KA-C217-bigi, KA-C217-sadila, KA-C217-dappa

\chapter{Cheap, Empty, Tight, Loose, Thick}
\label{ch:ka-cheap-empty-tight-loose-thick}
\hlchaptermodality{217}

\begin{chapteropening}
\textbf{By the end of this chapter:} \emph{I can say cheap, empty, tight, loose and thick.}

Complete the last lesson of chapter 217: I can say cheap, empty, tight, loose and thick.
\end{chapteropening}

\section[\texorpdfstring{\kn{ಅಗ್ಗ} (agga)}{ಅಗ್ಗ (agga)}]{\kn{ಅಗ್ಗ} (agga) — cheap}

\label{lesson:KA-C217-agga}



\begin{quote}
Before the new one: say the Kannada for responsibility, then the Kannada for a duty.
\end{quote}

\subsection*{You'll want to know: \kn{ಅಗ್ಗ}}
\textbf{\kn{ಅಗ್ಗ}} — \emph{agga} — \textquotedblleft{}cheap\textquotedblright{}.

\begin{grammarlens}[title={What you've built}]
One more word for this chapter.
\end{grammarlens}

\subsection*{Your turn}
\begin{itemize}
\raggedright
  \item \emph{Say it:} \emph{agga}
  \item \emph{Say it:} \emph{agga}, once more
  \item \emph{Say it:} say \emph{agga}
  \item \emph{Recall:} say the Kannada for responsibility, then the Kannada for a duty, then say \emph{agga} again
\end{itemize}

\begin{tcolorbox}[breakable,colback=teal!4,colframe=teal!35!black,title={Before you move on}]
What does \kn{ಅಗ್ಗ} mean? (\textquotedblleft{}Cheap\textquotedblright{}.) Say it once more.
\end{tcolorbox}
\section[\texorpdfstring{\kn{ಖಾಲಿ} (khāli)}{ಖಾಲಿ (khāli)}]{\kn{ಖಾಲಿ} (khāli) — empty}

\label{lesson:KA-C217-khali}



\begin{quote}
Before the new one: say the Kannada for a duty, then the Kannada for cheap.
\end{quote}

\subsection*{You'll want to know: \kn{ಖಾಲಿ}}
\textbf{\kn{ಖಾಲಿ}} — \emph{khāli} — \textquotedblleft{}empty\textquotedblright{}.

\begin{grammarlens}[title={What you've built}]
One more word for this chapter.
\end{grammarlens}

\subsection*{Your turn}
\begin{itemize}
\raggedright
  \item \emph{Say it:} \emph{khāli}
  \item \emph{Say it:} \emph{khāli}, once more
  \item \emph{Say it:} say \emph{khāli}
  \item \emph{Recall:} say the Kannada for a duty, then the Kannada for cheap, then say \emph{khāli} again
\end{itemize}

\begin{tcolorbox}[breakable,colback=teal!4,colframe=teal!35!black,title={Before you move on}]
What does \kn{ಖಾಲಿ} mean? (\textquotedblleft{}Empty\textquotedblright{}.) Say it once more.
\end{tcolorbox}
\section[\texorpdfstring{\kn{ಬಿಗಿ} (bigi)}{ಬಿಗಿ (bigi)}]{\kn{ಬಿಗಿ} (bigi) — tight}

\label{lesson:KA-C217-bigi}



\begin{quote}
Before the new one: say the Kannada for cheap, then the Kannada for empty.
\end{quote}

\subsection*{You'll want to know: \kn{ಬಿಗಿ}}
\textbf{\kn{ಬಿಗಿ}} — \emph{bigi} — \textquotedblleft{}tight\textquotedblright{}.

\begin{grammarlens}[title={What you've built}]
One more word for this chapter.
\end{grammarlens}

\subsection*{Your turn}
\begin{itemize}
\raggedright
  \item \emph{Say it:} \emph{bigi}
  \item \emph{Say it:} \emph{bigi}, once more
  \item \emph{Say it:} say \emph{bigi}
  \item \emph{Recall:} say the Kannada for cheap, then the Kannada for empty, then say \emph{bigi} again
\end{itemize}

\begin{tcolorbox}[breakable,colback=teal!4,colframe=teal!35!black,title={Before you move on}]
What does \kn{ಬಿಗಿ} mean? (\textquotedblleft{}Tight\textquotedblright{}.) Say it once more.
\end{tcolorbox}
\section[\texorpdfstring{\kn{ಸಡಿಲ} (saḍila)}{ಸಡಿಲ (saḍila)}]{\kn{ಸಡಿಲ} (saḍila) — loose}

\label{lesson:KA-C217-sadila}



\begin{quote}
Before the new one: say the Kannada for empty, then the Kannada for tight.
\end{quote}

\subsection*{You'll want to know: \kn{ಸಡಿಲ}}
\textbf{\kn{ಸಡಿಲ}} — \emph{saḍila} — \textquotedblleft{}loose\textquotedblright{}.

\begin{grammarlens}[title={What you've built}]
One more word for this chapter.
\end{grammarlens}

\subsection*{Your turn}
\begin{itemize}
\raggedright
  \item \emph{Say it:} \emph{saḍila}
  \item \emph{Say it:} \emph{saḍila}, once more
  \item \emph{Say it:} say \emph{saḍila}
  \item \emph{Recall:} say the Kannada for empty, then the Kannada for tight, then say \emph{saḍila} again
\end{itemize}

\begin{tcolorbox}[breakable,colback=teal!4,colframe=teal!35!black,title={Before you move on}]
What does \kn{ಸಡಿಲ} mean? (\textquotedblleft{}Loose\textquotedblright{}.) Say it once more.
\end{tcolorbox}
\section[\texorpdfstring{\kn{ದಪ್ಪ} (dappa)}{ದಪ್ಪ (dappa)}]{\kn{ದಪ್ಪ} (dappa) — thick, fat}

\label{lesson:KA-C217-dappa}



\begin{quote}
Before the new one: say the Kannada for tight, then the Kannada for loose.
\end{quote}

\subsection*{You'll want to know: \kn{ದಪ್ಪ}}
\textbf{\kn{ದಪ್ಪ}} — \emph{dappa} — \textquotedblleft{}thick, fat\textquotedblright{}.

\begin{grammarlens}[title={What you've built}]
One more word for this chapter.
\end{grammarlens}

\subsection*{Your turn}
\begin{itemize}
\raggedright
  \item \emph{Say it:} \emph{dappa}
  \item \emph{Say it:} \emph{dappa}, once more
  \item \emph{Say it:} say \emph{dappa}
  \item \emph{Recall:} say the Kannada for tight, then the Kannada for loose, then say \emph{dappa} again
\end{itemize}

\begin{tcolorbox}[breakable,colback=teal!4,colframe=teal!35!black,title={Before you move on}]
What does \kn{ದಪ್ಪ} mean? (\textquotedblleft{}Thick, fat\textquotedblright{}.) Say it once more.
\end{tcolorbox}
